\section{Introduction}\label{sec:intro}

Data centers are among the fastest-growing electricity consumers. Efficiency gains kept their energy use nearly flat through the 2010s even as computing output grew several-fold~\cite{masanet2020recalibrating,shehabi2016united}. That buffer is now thinning, and information and communication technology accounts for an estimated 2.1--3.9\% of global greenhouse-gas emissions~\cite{freitag2021real}. Large machine-learning workloads have drawn particular attention~\cite{strubell2019energy,patterson2022carbon,dodge2022measuring}. The carbon intensity of grid electricity varies by a factor of ten or more across regions and hours, driven by the output of wind, solar and dispatchable plant~\cite{staffell2017measuring}. Operators can therefore cut emissions without new hardware by moving flexible work to cleaner places and times. Such \emph{carbon-aware computing}~\cite{radovanovic2023carbon,bashir2021enabling} is an operations problem: capacity is finite, deadlines bind, and both the grid and the job stream are uncertain.

Batch jobs such as model training, analytics, rendering and backups are a natural target, because they tolerate some delay. Previous work has shown the potential of temporal shifting (running later)~\cite{wiesner2021lets,goiri2012greenhadoop,hanafy2023carbonscaler} and of spatial shifting (running elsewhere)~\cite{liu2011greening,zheng2020mitigating}. It has also shown their limits. Sukprasert et al.~\cite{sukprasert2024limitations} found that practical reductions from spatio-temporal shifting fall well short of idealized bounds. Most methods are evaluated with a scheduler that reacts to the jobs it has already seen, using point forecasts of carbon intensity. Three questions have received less attention. First, how much is lost because an online scheduler does not know which jobs will arrive next, and can any scheduler guarantee good performance in the worst case? Second, does it pay to price the uncertainty of carbon-intensity forecasts into scheduling decisions? Third, when several sites pool their capacity, how should the resulting emissions be attributed so that every participant is better off than on its own?

This paper addresses these questions with a scheduling policy, guarantees under explicitly stated assumptions, and a controlled evaluation on three years of hourly grid data for Great Britain, Europe, the United States, Brazil and New Zealand. Its building blocks (receding-horizon planning, network flows, conformal calibration and dual cost allocation) are standard; the contribution lies in their combination for this problem, the lower bound, the conditions under which the guarantees hold and the empirical evidence. It makes five contributions.
\begin{enumerate}
\item \textbf{Formulation.} We state multi-site, deadline-constrained, width-limited scheduling of deferrable bag-of-tasks jobs, including the emissions of data transfer, as a minimum-cost flow problem (Proposition~\ref{prop:flow}). Offline clairvoyant optima and every online planning step can therefore be solved to optimality, after quantization, by fast combinatorial min-cost flow algorithms. We also prove that myopic receding-horizon scheduling can be arbitrarily worse than the clairvoyant optimum even with perfect carbon forecasts (Proposition~\ref{prop:myopic}).
\item \textbf{Algorithm.} We propose CARMA (Carbon-aware Anticipatory Receding-horizon Multi-site Allocation). At every hour, CARMA adds \emph{ghost jobs} to its planning network. These represent expected arrivals, learned from the operator's history, so the planner reserves scarce low-carbon capacity for work that has not yet arrived. The planning problem stays a single min-cost flow and decisions take tens of milliseconds.
\item \textbf{Guarantees.} We prove that uncertainty about demand alone forces every online algorithm, even one with perfect carbon forecasts, to a competitive ratio of $\Omega(\sqrt{\theta})$, about $\sqrt{\theta}/2$ for large $\theta$, where $\theta$ is the ratio of the highest to the lowest unit emission (Theorem~\ref{thm:lb}). In an idealized setting with a backstop resource, CARMA never misses a deadline whatever its predictions, so it inherits the $\theta$-competitiveness that every on-time policy has. If, in addition, each episode fits within the planning window, CARMA is optimal when its predictions of demand and carbon are exact, and its excess emissions grow at most linearly in the prediction errors (Proposition~\ref{prop:robust} and Theorem~\ref{thm:carma}). These assumptions do not hold in our weekly experiments, so we verify the theorem separately in episodes that satisfy them. We also show that the dual prices of the pooled scheduling network give a carbon attribution among sites that lies in the core of the pooling game, so no group of sites gains by leaving (Theorem~\ref{thm:core}).
\item \textbf{Forecasting and recalibration pipeline.} As a tool for the scheduling study, we build a reproducible 1--24-hour carbon-intensity forecaster from a public feed: direct gradient-boosted models with horizon-wise normalized conformal quantiles and adaptive recalibration, which keeps empirical coverage close to target under the decarbonization drift of 2024. Its quantiles let us test risk-adjusted carbon pricing inside the scheduler. We do not claim a new forecasting method.
\item \textbf{Evidence design.} We evaluate CARMA out of sample on 52 weeks of 2024, with three workload traces per week, three load levels and two operating regimes. We add a trace-derived workload from the Alibaba 2018 cluster trace and four further fleets built from open generation data for Europe, the United States, Brazil and New Zealand. CARMA is compared with seven implementable online baselines, including a revenue-management protection-level policy and a scenario-based stochastic lookahead that uses the same demand information, with clairvoyant and ablated variants, and with the clairvoyant oracle. The experiments separate the value of anticipating demand from the value of better carbon forecasts, characterize how performance degrades with prediction errors, test the conclusions under a stylized marginal-emission accounting, and assess the stability of origin-based, host-based, Shapley and dual carbon attribution.
\end{enumerate}

The emission figures in this paper are attributional, based on average grid intensities, and Section~\ref{sec:sens} examines how they change under a stylized marginal-emission accounting. The rest of the paper is organized as follows. Section~\ref{sec:related} reviews related work. Section~\ref{sec:methods} presents the model, the forecaster, CARMA and the experimental design. Section~\ref{sec:theory} establishes the theoretical guarantees, with proofs in \appref{app:proofs}. Section~\ref{sec:results} reports results, Section~\ref{sec:discussion} discusses implications and limitations, and Section~\ref{sec:conclusion} concludes.

\section{Related work}\label{sec:related}

\paragraph{Carbon-aware workload shifting} Early green data-center schedulers matched deferrable jobs to on-site solar production, for example GreenSlot~\cite{goiri2011greenslot} and GreenHadoop~\cite{goiri2012greenhadoop}. Other work exploited electricity-price and carbon differences across sites through geographical load balancing~\cite{rao2010minimizing,liu2011greening} and dynamic right-sizing~\cite{lin2013dynamic}. With public carbon-intensity feeds, attention moved to grid emissions. Wiesner et al.~\cite{wiesner2021lets} quantified the reductions from temporal shifting in several regions, including Great Britain, and proposed forecast-based lowest-window scheduling. Google's carbon-intelligent compute system sets next-day virtual capacity curves from carbon forecasts and forecasts of the fleet's own demand~\cite{radovanovic2023carbon}. It is the closest industrial precedent for using demand forecasts, but it shapes hourly virtual capacity curves per cluster rather than planning individual jobs across sites. Cucumber admits delay-tolerant work only when forecast renewable supply allows it~\cite{wiesner2022cucumber}, and Caribou shifts serverless functions between regions~\cite{gsteiger2024caribou}. CarbonScaler varies the allocation of elastic jobs with carbon intensity~\cite{hanafy2023carbonscaler}. Ecovisor exposes the energy system to applications~\cite{souza2023ecovisor}. Carbon Explorer weighs scheduling against investment in renewables and storage~\cite{acun2023carbon}. Load migration between data centers can absorb curtailed renewable generation~\cite{zheng2020mitigating}, and data-center demand response has been surveyed in~\cite{wierman2014opportunities}. Sukprasert et al.~\cite{sukprasert2024limitations} bound the potential of spatio-temporal shifting across 123 regions and find that simple policies capture most of the achievable benefit. Our study complements theirs by identifying \emph{which} information limits online schedulers when capacity at clean sites is scarce.

\paragraph{Online algorithms and guarantees} Worst-case analysis of carbon-aware scheduling builds on online search and one-way trading~\cite{elyaniv2001optimal}, for which threshold policies achieve optimal competitive ratios. Lechowicz et al.\ derived optimal double-threshold algorithms for the online pause-and-resume problem~\cite{lechowicz2023online}. With CarbonClipper, they also derived learning-augmented algorithms with an optimal consistency--robustness trade-off for spatiotemporal allocation of a single deadline-constrained workload over a network of data centers~\cite{lechowicz2025learning}. These works model uncertainty about future carbon \emph{prices} for one job. We study many jobs that compete for finite low-carbon capacity, and we show that uncertainty about \emph{demand} creates a distinct lower bound that carbon forecasts cannot remove. Our guarantees follow the consistency--robustness paradigm of algorithms with predictions~\cite{lykouris2021competitive}. CARMA's re-solving of a deterministic plan with expected future demand is an instance of the certainty-equivalent re-solving heuristics of network revenue management~\cite{jasin2012resolving}. Our contribution is not the device itself but its formulation as a single flow problem for carbon-aware scheduling, the guarantees and the empirical attribution of value. Table~\ref{tab:related} positions the closest work along the dimensions that matter here.

\begin{table}[!htbp]
\centering
\caption{Positioning against the closest work, as modeled in the cited studies ($\checkmark$: explicitly modeled or provided; --: not modeled).}
\label{tab:related}
\footnotesize
\footnotesize\setlength{\tabcolsep}{3pt}
\begin{tabular}{p{4.3cm}cccc}
\toprule
Work & \makecell{Many jobs share\\finite capacity} & \makecell{Information on\\future demand} & \makecell{Worst-case\\guarantee} & \makecell{Attribution\\among sites} \\
\midrule
Lowest-window temporal shifting~\cite{wiesner2021lets} & -- & -- & -- & -- \\
Geographical load balancing~\cite{liu2011greening} & $\checkmark$ & -- & -- & -- \\
Carbon-intelligent compute system~\cite{radovanovic2023carbon} & $\checkmark$ & $\checkmark$ & -- & -- \\
Online pause-and-resume~\cite{lechowicz2023online} & -- & -- & $\checkmark$ & -- \\
CarbonClipper~\cite{lechowicz2025learning} & -- & -- & $\checkmark$ & -- \\
Flow models of preemptive scheduling~\cite{horn1974some,federgruen1986preemptive} & $\checkmark$ (offline) & -- & -- & -- \\
This paper & $\checkmark$ & $\checkmark$ & $\checkmark$ (A1--A2) & $\checkmark$ \\
\bottomrule
\end{tabular}
\end{table} Our fairness result draws on cooperative game theory: linear production games and flow games have non-empty cores that can be computed from dual prices~\cite{owen1975core,kalai1982totally}.

\paragraph{Energy-aware scheduling in operations} Energy-efficient scheduling is a mature topic in manufacturing operations~\cite{gahm2016energy}, and green scheduling models have appeared in this journal, for example for parallel machines with job splitting~\cite{asadpour2022green} and for cloud resource allocation~\cite{sharma2024efficient}. Most such models are deterministic and solved offline, often with metaheuristics. We instead exploit exact network-flow structure and embed it in an online lookahead policy. Network-flow formulations of preemptive scheduling with release dates and deadlines go back to Horn~\cite{horn1974some} and Federgruen and Groenevelt~\cite{federgruen1986preemptive}. Proposition~\ref{prop:flow} extends this classical construction to multiple sites, job widths, time-varying emission costs and transfer emissions~\cite{ahuja1993network,orlin1997polynomial}. In the taxonomy of Powell~\cite{powell2019unified}, CARMA is a deterministic lookahead policy with a modified (anticipatory) model, in the spirit of model predictive control~\cite{rawlings2017model,mayne2000constrained}.

\paragraph{Carbon-intensity forecasting} Forecasts of grid carbon intensity use statistical decomposition~\cite{bokde2021short}, machine learning with many exogenous drivers~\cite{leerbeck2020short}, day-ahead regression for building demand response~\cite{lowry2018day}, and hierarchical models of the generation mix, as in CarbonCast~\cite{maji2022carboncast}. For Great Britain, NESO publishes regional forecasts through its Carbon Intensity API~\cite{nesoapi}. These are average (attributional) intensities. Marginal emission factors, which describe the plant that responds to a change in load, can differ substantially from averages and are the relevant measure for the system-level effect of load shifting~\cite{hawkes2010estimating,silerevans2012marginal}. These studies focus on point accuracy. For decisions, calibrated uncertainty matters as well. Conformal prediction~\cite{vovk2005algorithmic,lei2018distribution,angelopoulos2023conformal} provides distribution-free intervals and has been extended to quantile regression~\cite{romano2019conformalized}, to multi-horizon~\cite{stankeviciute2021conformal} and dynamic time series~\cite{xu2023conformal}, and to arbitrary distribution shift through adaptive conformal inference~\cite{gibbs2021adaptive}. We use these tools to build a recalibrated forecaster and to test, in a controlled way, whether risk-adjusted carbon costs improve scheduling. The question is motivated by the optimizer's curse~\cite{smith2006optimizer} and by the robust-optimization view that point-estimate solutions can be fragile~\cite{bental2000robust,bertsimas2004price}.
